\paragraph{The same surface gate for all abelian schemes and all bases.}
The universal finiteness assertion in the manuscript is broader
than the relative-curve assertion just considered.  Its base may
be singular, nonreduced or nonseparated.  The following descent
argument shows that this added generality does not enlarge the
remaining divisor-Tate gate.

\begin{lemma}[Rational sections descend through proper covers]
\label{lem:B-proper-section-descent}
\status{Proved}
Let \(S\) be a noetherian \(\F_p\)-scheme, let
\(q:T\to S\) be proper and surjective, and let \(A/S\)
be an abelian scheme.  The sequence
\[
0\longrightarrow A(S)\otimes\Q
\longrightarrow A(T)\otimes\Q
\xrightarrow{\operatorname{pr}_1^*-\operatorname{pr}_2^*}
A(T\times_S T)\otimes\Q
\]
is exact.  The same equalizer statement holds after scalar
extension from \(\Q\) to \(\Q_\ell\).
\end{lemma}
\begin{proof}
First suppose \(S\) reduced, and take an actual section
\(b:T\to A\) whose two pullbacks agree as morphisms.
Let \(Z\) be its scheme-theoretic image in \(A\), with
maps \(j:Z\hookrightarrow A\) and \(g:Z\to S\).
The morphism \(g\) is proper and \(T\to Z\) is
surjective, by Stacks Lemma~29.42.10~\cite{BStacksProperImage}.
On every geometric fibre, any pair of points of \(T\) has
the same image under \(b\), by the equalizer hypothesis.
Since proper surjectivity persists under base change, each
geometric fibre of \(g\) has exactly one point.  Thus \(g\)
is quasi-finite and radicial as well as proper and surjective.
It is a finite universal homeomorphism; the exact criteria are
Stacks Lemmas~29.10.2 and~37.44.1~\cite{BStacksRadicial,BStacksProperFinite}.

There are an integer \(r\ge0\) and a morphism
\(h:S\to Z\) such that
\[
 h g=F_Z^r,\qquad g h=F_S^r.
\]
This is Kelly's Proposition~3 and Corollary~20, proved through
Lemma~18 and Corollary~19~\cite{BKelly2018}.
For reduced \(S\), it also follows directly from
Stacks Lemma~29.47.9: every local element of \(\mathcal O_Z\)
has a \(p\)-power in \(\mathcal O_S\), and finite
presentation and quasi-compactness give one exponent for a finite
set of algebra generators~\cite{BStacksFrobeniusPower}.
The local power maps glue.

Write
\(A^{(p^r)}=A\times_{S,F_S^r}S\).
The pair \((jh,1_S)\) defines a section \(c\) of
\(A^{(p^r)}/S\), and its pullback to \(T\) is
\(F_{A/S}^r b\).  Applying the iterated Verschiebung gives
\[
 a=V_{A/S}^r c\in A(S),\qquad q^*a=[p^r]b,
\]
because \(V^rF^r=[p^r]\); see Conrad's Example~4.4
and the Frobenius--Verschiebung discussion~\cite{BConrad2006}.
Thus an actual equalizer section descends after multiplication
by a power of \(p\).

For a rational equalizer element, multiply first by an integer
to obtain an actual section.  Its actual Cech difference is
torsion, since its rational difference vanishes.  Multiplication
by the order of that difference reduces to the preceding case.
Injectivity over reduced \(S\) follows because the equality
locus of two sections into the separated scheme \(A\) is
closed: surjective pullback makes its defining ideal nilpotent,
and reducedness makes that ideal zero.

For nonreduced \(S\), the finite filtration of its nilpotent
ideal has successive additive section kernels
\(\operatorname{Lie}(A)\otimes I\).  Its cohomology exact
sequences show that reduction changes section groups only by
bounded \(p\)-primary kernels and cokernels, as in
Theorem~\ref{thm:B-finite-chain-all-bases}.
Reduce \(T\) and its fibre product as well.  Rationalizing
therefore reduces the asserted sequence to the reduced case.
Finally \(\Q_\ell\) is flat over \(\Q\), so scalar
extension preserves the equalizer.  This last step avoids any
attempt to clear arbitrary \(\Q_\ell\)-coefficients by an
integer.
\end{proof}

\begin{theorem}[The universal all-base finiteness gate]
\label{thm:B-all-bases-surface-gate}
\status{Proved equivalence; the surface Tate assertion is unresolved}
Fix distinct primes \(p\) and \(\ell\).  The following
universal assertions are equivalent.
\begin{enumerate}
\item Divisor Tate holds for every smooth projective surface
over every finite extension of \(\F_p\).
\item For every finite-type \(\F_p\)-scheme \(S\) and
every abelian scheme \(A/S\), the group
\(H^1_{\et}(S,A)[\ell^\infty]\) is finite.  No
regularity, reducedness, properness or separatedness hypothesis
is imposed on \(S\).
\item The assertion in (2) holds for the Jacobians of smooth
projective geometrically connected curve families with a section
over nonempty opens of \(\mathbf P^1\), over all finite
extensions of \(\F_p\).
\end{enumerate}
These assertions are also equivalent to the birational
degree-three vanishing in
Theorem~\ref{thm:B-relative-surface-gate}.
\end{theorem}
\begin{proof}
The implication (2)\(\Rightarrow\)(3) is immediate.
Proposition~\ref{prop:B-relative-curve} and
Theorem~\ref{thm:B-relative-surface-gate} give
(3)\(\Rightarrow\)(1).  We prove (1)\(\Rightarrow\)(2).

First let \(T/k\) be smooth and separated of finite type,
where \(k/\F_p\) is finite, and let \(A/T\) be abelian.
Set \(B=A^t\), with structure map \(\pi:B\to T\).
The Kahn--Morrow implication used in
Theorem~\ref{thm:B-relative-surface-gate} gives divisor Tate
on the smooth separated total space \(B\).
For any integer \(m>1\), multiplication \([m]\) has
eigenvalue \(m^q\) on \(R^q\pi_*\Q_\ell\).
The distinct eigenvalues split the degree-two Leray filtration;
its \(m\)-eigenspace is
\[
 e_1 H^2(\bar B,\Q_\ell(1))
 =H^1(\bar T,R^1\pi_*\Q_\ell(1))
 =H^1(\bar T,V_\ell A).
\]
The divisor identification proved in
Proposition~\ref{prop:B-finite-chain-proper} identifies the
\(m\)-eigenspace of the algebraic divisor image with
\(N_T=A(\bar T)\otimes\Q_\ell\): the Poincare bundle
realizes sections, the zero section kills their Brauer obstruction,
and base Picard and relative Neron--Severi classes have eigenvalues
one and \(m^2\), respectively.  Applying \(e_1\) to
divisor Tate therefore proves
\[
 \kappa_T:N_T^\Gamma
 \xrightarrow{\sim}H^1(\bar T,V_\ell A)^\Gamma.
\]
Injectivity here is geometric Kummer, with the completion
identification in Lemma~\ref{lem:B-sections-reduced}.

Now let \(S/k\) be reduced, separated and of finite type.
Take a proper surjection \(q:T\to S\) with \(T\)
smooth separated of finite type: take the finitely many reduced
irreducible components, compactify, apply a projective alteration,
and restrict above \(S\).  De Jong's Theorem~4.1 and
Remark~4.2 supply the required smooth alterations~\cite{c_dejong1996}.
If a finite extension of constants is used, retain it in the
proper cover; assertion (1) includes all such finite extensions.
Put \(T_a=T^{\times_S(a+1)}\) and
\(\mathcal V=V_\ell A\).  Proper Cech descent gives
\[
 E_1^{a,b}=H^b(\bar T_a,\mathcal V)
 \ \Longrightarrow\ H^{a+b}(\bar S,\mathcal V).
\]
This follows from Stacks Lemmas~85.36.3--85.36.6 at finite
\(\ell\)-power coefficients~\cite{BStacksProperDescent}.
Finite-coefficient etale cohomology is finite in this situation,
so the inverse systems satisfy Mittag--Leffler.  Taking continuous
cohomology and tensoring with \(\Q_\ell\) preserves the
displayed first-quadrant spectral sequence in total degree one.
Its five-term sequence identifies the pullback kernel on
\(H^1\) with \(E_2^{1,0}\).

Each \(H^0(\bar T_a,\mathcal V)\) is pure of weight
\(-1\), including when \(T_a\) is singular: after a
finite constant extension, evaluation at one closed point of
each geometric connected component embeds it equivariantly into
a finite sum of Tate modules of abelian varieties over finite
fields.  Consequently \(E_2^{1,0}\), a subquotient of the
degree-zero row, has no \(\Gamma\)-invariants.

Take \(c\in H^1(\bar S,\mathcal V)^\Gamma\).
The smooth-base case gives a unique \(b\in N_T^\Gamma\)
such that \(q^*c=\kappa_T(b)\).
The two pullbacks of \(c\) to \(\bar T_1\) agree.
Kummer injectivity on the possibly singular \(T_1\), supplied
by Theorem~\ref{thm:B-finite-chain-all-bases}, therefore
gives \(\operatorname{pr}_1^*b=\operatorname{pr}_2^*b\).
Lemma~\ref{lem:B-proper-section-descent}, applied over
\(\bar k\) and then extended from \(\Q\) to
\(\Q_\ell\), descends \(b\) uniquely to
\(a\in N_S\).  Uniqueness makes \(a\) invariant.
The difference \(c-\kappa_S(a)\) is then invariant in
the weight-minus-one pullback kernel, hence zero.  Thus
invariant Kummer surjectivity holds over every reduced separated
base.

Nilpotent thickening does not change finite etale torsion sheaves
or their cohomology, and changes section groups only by bounded
\(p\)-primary groups.  This removes reducedness.  For a
nonseparated finite-type \(S\), use a finite affine open
cover.  Every finite intersection is quasi-compact and separated.
The ordinary open-cover Cech spectral sequence again has its
degree-one pullback kernel in the pure weight-minus-one
degree-zero row.  The pulled-back section descends because
sections are a Zariski sheaf and tensoring with \(\Q\)
is exact for the finite products in its equalizer, followed by
scalar extension to \(\Q_\ell\).
The same subtraction argument proves invariant Kummer
surjectivity on \(S\).

Finally Theorem~\ref{thm:B-finite-chain-all-bases}, including
its finite-cotype assertion, identifies this invariant Kummer
surjectivity with the actual finiteness of
\(H^1(S,A)[\ell^\infty]\).  This proves (2).
\end{proof}

This is a reduction of the full finite-type abelian-scheme
finiteness assertion, not a proof of its equivalent surface Tate
condition.  Extending the relative-curve Brauer quotient
statement to singular bases requires separate control of
non-torsion Picard-degree and base \(H^2(S,\Gm)\) terms.
The integral torsor-lattice and unit-cup arguments below provide
that control and prove the same surface-Tate equivalence in
the full arbitrary-base scope.  The characteristic-primary
counterexample below is separate.
